\section{Simulation Methods}
\label{sec:methods}

\subsection{Mode solver and geometry}

Modes were computed with femwell~\cite{femwell}, a full-vector finite-element mode solver built on
scikit-fem~\cite{gustafsson2020}, on conforming triangular meshes generated by Gmsh~\cite{geuzaine2009}. The mesh
follows every material boundary, so a perturbed wall is represented exactly, without a staircase. The
production mesh (``r12'') has a \qty{12}{nm} element size in the core. The domain spans
$|x|\le\qty{1.4}{\micro m}$ and $|y|\le\qty{1.2}{\micro m}$, with silicon $\nsi=3.48$ and oxide $\nox=1.44$ for both the buried oxide and the
cladding, without material dispersion, at $\lambda=\qty{1.55}{\micro m}$. These indices are rounded assumptions; the
measured bulk values at \qty{1.5}{\micro m} and \qty{295}{K} are \num{3.48314} for silicon~\cite{frey2006} and
\num{1.44502} for fused silica~\cite{leviton2008}. The width $w$ is the bottom width of the wall; the angle tilts the
wall about its bottom edge, which for a rib is the slab top. Walls are ideal trapezoids. For each geometry we solve for
two modes near a guess and keep the highest-index mode of the requested polarization, falling back to a six-mode search
without a guess if that mode is more than \qty{3}{\percent} from the guess. The group index is
$\ngr=\neff-\lambda\,\partial\neff/\partial\lambda$ from a central difference over \qty{\pm 5}{nm} on the same mesh.

\subsection{Derivatives and their checks}

The Jacobian in~\eqref{eq:fisher} is built from central differences, each perturbed geometry being remeshed, with
steps $h$ of \qty{5}{nm} in width, \qty{2}{nm} in thickness and slab, \qty{1}{\degree} in angle and $0.002$ in each
index. Every derivative was also computed with steps $h/2$ and $2h$. Of the 428 derivatives, 422 agree to within
\qty{2}{\percent} over the three steps, 3 more are near zero and agree in absolute terms, and 3 exceed the limit
slightly, the largest by \qty{2.25}{\percent}; all three are $\ngr$ derivatives. Bounds are reported at step $h$ with
the range over $h/2$, $h$ and $2h$.

Refining the mesh from r12 to r8 (\qty{8}{nm} core elements) changes no derivative of the strip and rib TE
\qty{500}{nm} structures by more than \qty{0.31}{\percent}. Because remeshing every perturbed geometry adds solve-to-solve
noise, we also recomputed the geometric derivatives of the \qty{300}{nm} ribs, the most weakly confined structures,
on a fixed mesh. The nominal mesh is reused, and its nodes are moved onto each perturbed geometry by a continuous
piecewise-linear map that keeps every material boundary on its line. The fixed-mesh derivatives differ from the
remeshed ones by at most \qty{0.13}{\percent} on the \qty{50}{nm} slab and \qty{0.78}{\percent} on the \qty{90}{nm}
slab. Enlarging the computational domain by a factor of 1.5 in each direction, to
$|x|\le\qty{2.1}{\micro m}$ and $|y|\le\qty{1.8}{\micro m}$, changes $\neff$ of the seven structures of the two-width set by
at most \num{2.04e-6} and $\ngr$ by at most \num{1.32e-5} in relative terms, and their geometric derivatives by at most
\qty{0.35}{\percent}. Because the two-width bound responds to derivative errors of this size, it moves by
\qty{5.2}{\percent}, to \qty{0.409}{\degree}, when the large-domain geometric derivatives replace the standard ones (the
index derivatives are kept from the standard domain).

The mode indices agree with an earlier finite-difference solution of the same structures (EMpy~\cite{empy}, with
cell-exact averaging): within \qty{0.79}{\percent} in $\neff$ for the ten strips and \qty{0.15}{\percent} for the five ribs.
For the rib of \qty{500}{nm} width on a \qty{90}{nm} slab, femwell gives $\neff=2.5707$, \qty{1.2}{\percent} below the
effective-index-method value of 2.602 reported for the same geometry and indices~\cite{jiao2025}; a full-vector solution
is expected to lie below the effective-index estimate.

\subsection{Why the strip symmetry is imposed}

Agreement in $\neff$ does not imply agreement in derivatives, and the analysis depends on derivatives. Across 30
structures, the EMpy derivatives differ from the femwell ones by a median of \qtyrange{4.7}{14}{\percent} depending on the
derivative, and in the same \qty{2}{\percent} step-size test 222 EMpy derivatives failed (373 passed and 53 were near
zero). We therefore use
femwell throughout. Even femwell's derivatives are not exact. The femwell strip ratios
$(\partial\neff/\partial\theta)/(\partial\neff/\partial w)$ scatter by up to \qty{0.84}{\percent} around the exact
value. That scatter alone is enough to create apparent angle information in strips. In an earlier pass of this
analysis, with placeholder noise values, the best set of strips alone gave a finite \qty{0.6135}{\degree} from the raw
femwell derivatives and an infinite bound once the exact relation was imposed. Since the relation is exact for oxide-clad
strips (Section~\ref{sec:theory}), every result below replaces each strip's angle column by $-t\tan(\qty{1}{\degree})$
times its width column.

\subsection{Monte Carlo fits}

To test whether the linearized bound is reached by an actual estimator, we fitted simulated noisy data with the full
nonlinear model. Because thousands of fits need fast forward evaluations, each structure's $\neff$ and $\ngr$ were
represented by a polynomial in its own inputs (width, thickness, angle, slab, $\nsi$, $\nox$ for ribs; mid-height width,
thickness and indices for strips, which encodes the symmetry), fitted to femwell solves on a Latin hypercube in a box of
$\pm6$ marginal CRLB standard deviations around nominal, with the polynomial degree chosen by leave-one-out error. On
held-out solves the strip surrogates are within $0.33\,\sigma_\phi$ in $\neff$ (with $\sigma_\phi$ expressed as an
$\neff$ noise) and the 400--\qty{600}{nm} ribs within $1.13$--$2.24\,\sigma_\phi$. The \qty{300}{nm} ribs miss by
$7.23\,\sigma_\phi$ (\qty{90}{nm} slab) and $27\,\sigma_\phi$ (\qty{50}{nm} slab). For these weakly confined ribs the
misses come from femwell's remeshing scatter between unrelated geometries, which no smooth model can follow; the same
scatter cancels within one finite-difference stencil, as the fixed-mesh check shows. In this Monte Carlo, data and fits
both use the surrogate, so it tests the estimator on a smoothed model; Section~\ref{sec:surrcheck} repeats the test
with data taken from direct femwell solves. Each case uses 2000 noise realizations at
$\sigma_\lambda=\qty{1}{pm}$, with prior centers drawn from the priors, and bounded nonlinear least squares from five
starts (nominal and four random points in the box). A fit is flagged when it ends at the box edge, and as ambiguous when
another start reaches a $\chi^2$ within 1 of the best with a different angle.
